% 图 49.1 —— 由 `checks/emit_hand.py` 自动拼出（probe 精测 + prims2 图元分解）
%%
% ⚠ 这是**稿子不是成品**：跑 `checks/checkhand.py 49.1` 看指标 + 看叠合图。
%%
%% ⚠ 待核：
%   · 本图 probe 报出阴影族 1 个 —— **本脚本不生成阴影**（要照原书手写）；prims2 的 strokes 里可能含阴影线，核对时留意别画重
%   · 可疑字形 1 项（空心圈 ⊙ 常被认成字母 o）—— 见文件末
%%
\documentclass[12pt]{standalone}
\usepackage{fontspec}
\setsansfont{Arial}
\newfontfamily\mfallback{DejaVu Sans}
\usepackage{tikz}
\begin{document}
\begin{tikzpicture}[x=1pt, y=-1pt, line width=4.0pt, line cap=round, line join=round]

  %% ════ 填充区（prims2 的 fills）════

  %% ════ 直线段（prims2 的 strokes）════
  \draw[line width=4.0pt] (191.0,403.0) -- (191.0,377.0);
  \draw[line width=6.0pt] (190.0,404.0) -- (132.0,405.0);
  \draw[line width=4.0pt] (190.0,364.0) -- (191.0,340.0);
  \draw[line width=4.0pt] (327.0,404.0) -- (275.0,404.0);
  \draw[line width=10.0pt] (196.0,93.0) -- (206.0,116.0);
  \draw[line width=5.0pt] (838.0,310.0) -- (852.0,309.0);
  \draw[line width=5.0pt] (130.0,242.0) -- (130.0,265.0);
  \draw[line width=5.0pt] (49.0,371.0) -- (44.0,371.0);
  \draw[line width=5.0pt] (133.0,66.0) -- (189.0,77.0);
  \draw[line width=4.0pt] (129.0,71.0) -- (129.0,67.0);
  \draw[line width=5.0pt] (1122.0,402.0) -- (1060.0,402.0);
  \draw[line width=4.0pt] (231.0,204.0) -- (205.0,118.0);
  \draw[line width=6.0pt] (131.0,124.0) -- (130.0,102.0);
  \draw[line width=6.0pt] (946.0,397.0) -- (893.4,320.4);
  \draw[line width=3.9pt] (893.4,320.4) -- (890.0,319.0);
  \draw[line width=8.2pt] (46.0,379.0) -- (43.0,372.0);
  \draw[line width=5.0pt] (548.0,402.0) -- (547.0,388.0);
  \draw[line width=5.0pt] (190.0,196.0) -- (191.0,216.0);
  \draw[line width=4.5pt] (131.0,137.0) -- (131.0,159.0);
  \draw[line width=5.8pt] (179.0,54.0) -- (173.8,51.8);
  \draw[line width=6.0pt] (722.0,404.0) -- (662.0,403.0);
  \draw[line width=3.0pt] (131.0,368.0) -- (131.0,347.0);
  \draw[line width=7.0pt] (180.0,55.0) -- (190.0,76.0);
  \draw[line width=5.0pt] (191.0,254.0) -- (190.0,233.0);
  \draw[line width=6.0pt] (948.0,397.0) -- (1000.0,319.0);
  \draw[line width=5.0pt] (191.0,146.0) -- (191.0,124.0);
  \draw[line width=4.0pt] (1002.0,386.0) -- (1002.0,402.0);
  \draw[line width=5.5pt] (946.0,400.0) -- (940.1,402.9);
  \draw[line width=5.0pt] (940.1,402.9) -- (889.0,402.0);
  \draw[line width=4.0pt] (191.0,161.0) -- (191.0,183.0);
  \draw[line width=5.0pt] (441.0,394.0) -- (440.0,219.0);
  \draw[line width=5.0pt] (1057.0,401.0) -- (1002.0,402.0);
  \draw[line width=3.0pt] (131.0,171.0) -- (131.0,194.0);
  \draw[line width=10.3pt] (195.0,92.0) -- (194.0,82.0);
  \draw[line width=5.0pt] (190.0,269.0) -- (191.0,290.0);
  \draw[line width=4.0pt] (547.0,403.0) -- (444.2,404.0);
  \draw[line width=9.0pt] (444.2,404.0) -- (442.0,395.0);
  \draw[line width=7.0pt] (50.0,372.0) -- (47.0,379.0);
  \draw[line width=5.0pt] (130.0,334.0) -- (132.0,312.0);
  \draw[line width=5.0pt] (1001.0,313.0) -- (1001.0,318.0);
  \draw[line width=5.0pt] (191.0,325.0) -- (190.0,305.0);
  \draw[line width=5.0pt] (130.0,228.0) -- (130.0,207.0);
  \draw[line width=5.0pt] (457.0,217.0) -- (441.0,218.0);
  \draw[line width=5.0pt] (43.0,370.0) -- (45.0,43.0);
  \draw[line width=3.0pt] (429.0,219.0) -- (412.0,219.0);
  \draw[line width=4.0pt] (94.0,58.0) -- (128.0,65.0);
  \draw[line width=5.0pt] (132.0,277.0) -- (131.0,298.0);
  \draw[line width=5.0pt] (129.0,72.0) -- (131.0,89.0);
  \draw[line width=5.0pt] (191.0,404.0) -- (248.0,404.0);
  \draw[line width=5.0pt] (549.0,403.0) -- (660.0,403.0);
  \draw[line width=6.0pt] (239.0,215.0) -- (237.0,203.0);
  \draw[line width=6.0pt] (888.0,319.0) -- (840.0,394.0);
  \draw[line width=5.3pt] (815.0,323.0) -- (816.0,318.0);
  \draw[line width=10.0pt] (262.0,317.0) -- (252.0,290.0);
  \draw[line width=8.0pt] (252.0,290.0) -- (239.0,217.0);
  \draw[line width=5.0pt] (889.0,386.0) -- (889.0,402.0);
  \draw[line width=5.0pt] (249.0,403.0) -- (250.0,380.0);
  \draw[line width=6.3pt] (550.0,217.0) -- (549.0,223.0);
  \draw[line width=4.0pt] (245.0,89.0) -- (195.0,81.0);
  \draw[line width=3.4pt] (232.0,204.0) -- (235.0,203.0);
  \draw[line width=5.0pt] (837.0,309.0) -- (836.0,18.0);
  \draw[line width=5.0pt] (231.0,205.0) -- (238.0,216.0);
  \draw[line width=5.0pt] (441.0,18.0) -- (440.0,217.0);
  \draw[line width=4.0pt] (250.0,404.0) -- (273.0,404.0);
  \draw[line width=6.0pt] (1002.0,319.0) -- (1057.0,400.0);
  \draw[line width=4.0pt] (45.0,43.0) -- (45.0,18.0);
  \draw[line width=5.0pt] (45.0,43.0) -- (62.0,42.0);
  \draw[line width=6.0pt] (548.0,224.0) -- (447.0,391.8);
  \draw[line width=4.8pt] (447.0,391.8) -- (443.0,394.0);
  \draw[line width=4.0pt] (837.0,310.0) -- (836.0,391.8);
  \draw[line width=3.0pt] (836.0,391.8) -- (838.0,394.0);
  \draw[line width=3.0pt] (131.0,381.0) -- (131.0,404.0);
  \draw[line width=5.0pt] (178.0,37.0) -- (180.0,54.0);
  \draw[line width=5.9pt] (207.0,117.0) -- (211.8,119.8);
  \draw[line width=6.0pt] (211.8,119.8) -- (236.0,202.0);
  \draw[line width=5.0pt] (550.0,224.0) -- (553.8,226.8);
  \draw[line width=6.0pt] (553.8,226.8) -- (656.0,390.0);
  \draw[line width=6.0pt] (656.0,390.0) -- (661.0,402.0);
  \draw[line width=4.2pt] (194.0,80.0) -- (191.0,77.0);
  \draw[line width=6.0pt] (274.0,403.0) -- (261.0,319.0);
  \draw[line width=4.9pt] (124.2,74.8) -- (128.0,72.0);
  \draw[line width=8.0pt] (46.0,380.0) -- (45.5,403.5);
  \draw[line width=4.0pt] (45.5,403.5) -- (130.0,405.0);
  \draw[line width=4.0pt] (250.0,334.0) -- (249.0,310.0);
  \draw[line width=6.0pt] (839.0,395.0) -- (839.0,401.0);
  \draw[line width=5.0pt] (839.0,401.0) -- (889.0,402.0);
  \draw[line width=6.0pt] (191.0,107.0) -- (194.0,93.0);
  \draw[line width=4.0pt] (250.0,368.0) -- (249.0,344.0);
  \draw[line width=6.0pt] (889.0,312.0) -- (889.0,318.0);

  %% ════ 曲线（prims2 的 curves —— 三段式，坐标同为 (y,x)）════
  \draw[line width=6.0pt] (173.8,51.8) .. controls (160.8,37.1) and (137.4,47.1) .. (132.0,64.0);
  \draw[line width=5.0pt] (134.0,419.0) .. controls (139.6,417.6) and (140.0,424.8) .. (138.0,428.0);
  \draw[line width=4.0pt] (271.0,336.0) .. controls (269.2,329.8) and (269.3,321.4) .. (263.0,318.0);
  \draw[line width=5.0pt] (1002.0,402.0) .. controls (984.6,402.6) and (963.1,406.0) .. (947.0,400.0);
  \draw[line width=6.0pt] (50.0,371.0) .. controls (71.6,271.3) and (87.2,169.9) .. (124.2,74.8);

  %% ════ 实心点 ════
  \fill (814.5,325.5) circle (5.7);
  \fill (250.0,407.5) circle (6.0);
  \fill (274.5,407.0) circle (5.5);
  \fill (662.5,407.0) circle (4.5);
  \fill (131.0,408.5) circle (5.0);
  \fill (191.0,408.0) circle (5.0);

  %% ════ 标签（probe 直接给的 \node 行）════
  %% ⚠ 全部补了 fill=white：文字在最上层，否则与曲线交叉干扰
     \node[anchor=center,inner sep=0pt,fill=white,font=\fontsize{32.6}{40.8}\selectfont] at (26.0,45.5) {\textsf{\textbf{\textit{α}}}};   % box(17,36,19,18)
     \node[anchor=center,inner sep=0pt,fill=white,font=\fontsize{27.8}{34.7}\selectfont] at (248.5,150.0) {\textsf{\textbf{\textit{f}}}};   % box(243,138,10,23)
     \node[anchor=center,inner sep=0pt,fill=white,font=\fontsize{27.2}{33.9}\selectfont] at (424.9,209.4) {\textsf{\textbf{\textit{a}}}};   % box(417,201,14,14)
     \node[anchor=center,inner sep=0pt,fill=white,font=\fontsize{23.6}{29.6}\selectfont] at (423.1,232.9) {\textsf{\textbf{2}}};   % box(417,224,11,17)
     \node[anchor=center,inner sep=0pt,fill=white,font=\fontsize{34.8}{43.5}\selectfont] at (628.0,302.4) {\textsf{\textbf{9}}};   % box(617,290,18,23)
     \node[anchor=center,inner sep=0pt,fill=white,font=\fontsize{31.0}{38.8}\selectfont] at (1026.8,304.1) {\textsf{\textbf{\textit{k}}}};   % box(1018,291,17,23)
     \node[anchor=center,inner sep=0pt,fill=white,font=\fontsize{29.1}{36.4}\selectfont] at (818.5,301.0) {\textsf{\textbf{\textit{a}}}};   % box(810,292,15,15)
     \node[anchor=center,inner sep=0pt,fill=white,font=\fontsize{31.2}{39.0}\selectfont] at (822.5,315.7) {{\mfallback\textbf{−}}};   % box(806,309,18,4)
     \node[anchor=center,inner sep=0pt,fill=white,font=\fontsize{21.5}{26.9}\selectfont] at (283.0,424.9) {\textsf{\textbf{1}}};   % box(277,417,7,15)
     \node[anchor=center,inner sep=0pt,fill=white,font=\fontsize{28.9}{36.1}\selectfont] at (664.9,433.0) {\textsf{\textbf{1}}};   % box(657,422,9,21)
     \node[anchor=center,inner sep=0pt,fill=white,font=\fontsize{29.7}{37.1}\selectfont] at (1064.5,433.5) {\textsf{\textbf{1}}};   % box(1056,423,10,20)
     \node[anchor=center,inner sep=0pt,fill=white,font=\fontsize{36.0}{45.0}\selectfont] at (258.1,439.1) {\textsf{\textbf{+}}};   % box(247,428,18,18)
     \node[anchor=center,inner sep=0pt,fill=white,font=\fontsize{30.1}{37.6}\selectfont] at (88.5,439.5) {\textsf{\textbf{\textit{x}}}};   % box(79,429,17,17)
     \node[anchor=center,inner sep=0pt,fill=white,font=\fontsize{30.1}{37.6}\selectfont] at (233.5,439.5) {\textsf{\textbf{\textit{x}}}};   % box(224,429,17,17)
     \node[anchor=center,inner sep=0pt,fill=white,font=\fontsize{30.1}{37.6}\selectfont] at (189.5,440.5) {\textsf{\textbf{\textit{x}}}};   % box(180,430,17,17)
     \node[anchor=center,inner sep=0pt,fill=white,font=\fontsize{31.2}{39.0}\selectfont] at (119.5,441.7) {{\mfallback\textbf{−}}};   % box(103,435,18,4)
     \node[anchor=center,inner sep=0pt,fill=white,font=\fontsize{31.2}{39.0}\selectfont] at (146.5,441.7) {{\mfallback\textbf{−}}};   % box(130,435,18,4)
     \node[anchor=center,inner sep=0pt,fill=white,font=\fontsize{31.2}{39.0}\selectfont] at (288.5,441.7) {{\mfallback\textbf{−}}};   % box(272,435,18,4)
     \node[anchor=center,inner sep=0pt,fill=white,font=\fontsize{22.7}{28.3}\selectfont] at (282.1,449.6) {\textsf{\textbf{6}}};   % box(276,441,11,16)
     \node[anchor=center,inner sep=0pt,fill=white,font=\fontsize{19.2}{24.0}\selectfont] at (140.3,450.3) {\textsf{\textbf{\textit{A}}}};   % box(133,442,11,16)

  % 钉死画布：否则超出的图元/控制点会把页面撑大，1:1 回比就失准
  \pgfresetboundingbox
  \path[use as bounding box] (0,0) rectangle (1137,463);
\end{tikzpicture}
\end{document}

% ⚠ 待核清单（probe 拿不准，人眼过一遍）：
%   · 未识别/跳过：⑤ 跳过/未识别的字形
